% ================================================
% 文件: 05_集成电路.tex
% 章节: 集成电路
% 所属部分: 电子元器件与材料
% 版本: v1.0
% ================================================

\section{集成电路}
\label{sec:integrated_circuits}

集成电路（IC）是将大量电子元件集成在一小块半导体芯片上。
所谓“集成”，是指采用平面工艺在同一片半导体基片上同时制作晶体管、
二极管、电阻、电容以及连接它们的金属互连线，使一个完整的电路单元
在物理上不可分割地形成于芯片内部，再经封装引出若干管脚构成一个
独立器件。与用分立元件搭接的电路相比，集成电路的体积和重量可缩小
几个数量级，功耗大幅降低，而可靠性显著提高——因为绝大部分故障都
源于焊点与引线，而集成电路把成千上万个连接点转移到了芯片内部的
金属化图形上。

\textbf{发展脉络}：集成电路的思想在二十世纪五十年代末提出，
其后半个世纪的发展基本遵循摩尔定律，即单块芯片上可容纳的晶体管
数目大致每十八至二十四个月翻一番。若以 $T$ 表示翻番周期，
则芯片上的元件数随时间近似满足
\[
N(t) = N_0 \cdot 2^{\,t/T}
\]
这一指数增长的物理基础是特征尺寸的持续缩小：工艺线宽从早期的
数十微米逐步缩小到微米、亚微米乃至纳米量级。线宽每缩小到原来的
$1/k$，同样面积上可容纳的器件数便增加约 $k^2$ 倍，同时门延迟按
比例减小、单个门的开关功耗按比例下降。正是这种“缩小即改善”的
标度律，使集成电路在性能、成本、功耗三方面同时获益，从而彻底
改变了通信设备的形态：昔日占满机柜的整机功能，如今往往集中在
一两枚芯片之中。

\textbf{按集成度分级}：集成度是衡量集成电路复杂程度的首要指标，
通常以单块芯片所含的元件数或等效逻辑门数计。按集成度由低到高，
习惯上作如下划分：

\begin{itemize}
    \item \textbf{SSI}：小规模集成电路（<100个元件）
    \item \textbf{MSI}：中规模集成电路（100-1000个元件）
    \item \textbf{LSI}：大规模集成电路（1000-100000个元件）
    \item \textbf{VLSI}：超大规模集成电路（>100000个元件）
\end{itemize}

小规模集成电路（SSI）所含元件不足一百个，典型产品是基本逻辑门、
触发器和单运算放大器，其功能相当于分立元件电路的直接替代，仍需
用户自行组合成系统。中规模集成电路（MSI）含一百至一千个元件，
已能实现计数器、译码器、多路开关、寄存器等具有独立功能的部件。
大规模集成电路（LSI）含一千至十万个元件，可把整个运算器、存储器
阵列或专用信号处理通道纳入单片。超大规模集成电路（VLSI）的元件
数超过十万，微处理器、大容量存储器与复杂的通信基带芯片均属此列；
其后又衍生出特大规模（ULSI）与巨大规模（GSI）等称谓，用以描述
元件数达到千万乃至十亿量级的芯片。需要指出，这种分级只是工程上
的粗略约定，随着工艺进步，各档的界线并无严格的物理意义，但它
清晰地反映了集成技术由“替代元件”到“替代整机”的演进轨迹。

\begin{table}[htbp]
\centering
\caption{集成电路规模分级与典型产品}
\begin{tabular}{|p{2.0cm}|p{3.0cm}|p{2.6cm}|p{4.0cm}|}
\hline
\textbf{级别} & \textbf{中文名称} & \textbf{元件数} & \textbf{典型产品} \\
\hline
SSI & 小规模集成电路 & 少于 100 & 基本逻辑门、触发器、单运放 \\
\hline
MSI & 中规模集成电路 & 100$\sim$1000 & 计数器、译码器、多路开关、寄存器 \\
\hline
LSI & 大规模集成电路 & 1000$\sim$100000 & 小型存储器、运算单元、专用通道 \\
\hline
VLSI & 超大规模集成电路 & 多于 100000 & 微处理器、大容量存储器、基带芯片 \\
\hline
ULSI/GSI & 特大规模、巨大规模 & 千万至十亿量级 & 片上系统、高性能处理器、FPGA \\
\hline
\end{tabular}
\end{table}

\textbf{按功能分类}：除按集成度划分外，更实用的分类方法是按处理
信号的性质划分。模拟集成电路处理连续变化的信号，输出与输入之间
保持确定的函数关系，代表性产品有运算放大器、电压比较器、稳压器、
音频功率放大器、混频器与锁相环。数字集成电路处理离散的二值信号，
只关心高、低两种电平，代表性产品有门电路、触发器、计数器、
存储器与微处理器。介于两者之间的是数模混合集成电路，其核心是
模数转换器（ADC）与数模转换器（DAC），负责在模拟世界与数字世界
之间建立桥梁。对于理想的 $N$ 位均匀量化转换器，其量化信噪比为
\begin{equation}
\mathrm{SNR} \approx 6.02N + 1.76 \ \ (\mathrm{dB})
\end{equation}
即每增加一位分辨力，信噪比约提高 $6\ \mathrm{dB}$。同时，为使
采样后能够无失真恢复原信号，采样频率须满足奈奎斯特条件
$f_s \geq 2f_{\max}$，工程实践中通常取更高的采样率并在前端加入
抗混叠低通滤波器。

\textbf{按工艺分类}：从制造工艺看，集成电路主要分为双极型、
金属氧化物半导体型与两者结合的 BiCMOS 型。双极型工艺以双极
晶体管为基本单元，速度快、驱动能力强、噪声低、模拟精度高，
但功耗大、集成度受限，早期的 TTL 逻辑电路与至今仍广泛使用的
高性能运算放大器多采用此工艺。MOS 工艺又分为 NMOS、PMOS 和
互补型 CMOS。CMOS 电路由 N 沟道与 P 沟道 MOSFET 成对构成，
静态时两管中总有一管截止，故静态功耗极小，其动态功耗主要来自
负载电容的充放电与短路电流，可近似表示为
\begin{equation}
P = \alpha C_L V_{DD}^2 f + I_{leak} V_{DD}
\end{equation}
式中 $\alpha$ 为翻转活动因子，$C_L$ 为等效负载电容，$V_{DD}$ 为
电源电压，$f$ 为工作频率。由该式可见，降低电源电压对节省功耗
最为有效（呈平方关系），这正是集成电路工作电压由 $5\ \mathrm{V}$
逐步降至 $3.3\ \mathrm{V}$、$1.8\ \mathrm{V}$ 乃至更低的原因；
同时也说明高速与低功耗之间存在固有矛盾。BiCMOS 工艺在同一芯片
上兼容双极管与 CMOS 管，用双极部分实现高速模拟与强驱动，用
CMOS 部分实现低功耗逻辑，在射频收发芯片中应用甚广。此外，
以砷化镓、氮化镓为衬底的化合物半导体集成电路，因电子迁移率高、
耐压大，专用于微波单片集成电路（MMIC）与大功率射频放大。

\begin{table}[htbp]
\centering
\caption{主要集成电路工艺的比较}
\begin{tabular}{|p{2.2cm}|p{2.6cm}|p{2.4cm}|p{4.2cm}|}
\hline
\textbf{工艺} & \textbf{静态功耗} & \textbf{速度/驱动} & \textbf{主要特点与用途} \\
\hline
双极型（TTL） & 大 & 快，驱动强 & 抗干扰强，模拟精度高；用于运放、高速逻辑 \\
\hline
CMOS & 极小 & 中等至很快 & 集成度高，电源范围宽，怕静电；数字电路主流 \\
\hline
BiCMOS & 中等 & 快 & 兼具双极高速与 CMOS 低功耗；射频收发、混合信号 \\
\hline
化合物半导体 & 较大 & 极快 & 迁移率高、耐压高；微波单片集成电路与功放 \\
\hline
\end{tabular}
\end{table}

\textbf{封装形式与散热}：芯片本身极为脆弱，必须经过封装才能使用。
封装的作用有四：提供机械保护、引出电气连接、传导内部热量、
隔绝湿气与污染。按安装方式可分为插装式与表面贴装式。插装式
封装以双列直插（DIP）为代表，管脚穿过印制板通孔焊接，便于
更换与调试，但占用面积大。表面贴装封装包括小外形封装（SOP）、
四边扁平封装（QFP）、球栅阵列（BGA）与芯片尺寸封装（CSP）等，
管脚间距小、寄生参数低、适于自动化装配，是现代设备的主流。
封装的热性能用热阻 $R_{th}$ 描述，芯片结温为
\[
T_j = T_a + P \cdot R_{th}
\]
其中 $T_a$ 为环境温度，$P$ 为耗散功率。集成电路的结温上限通常为
$125\ ^{\circ}\mathrm{C}$ 至 $150\ ^{\circ}\mathrm{C}$，而工作
温度每升高 $10\ ^{\circ}\mathrm{C}$，失效率大致翻倍，因此功率
较大的芯片必须加装散热片，或通过印制板上的散热焊盘与过孔阵列
把热量导出。

\textbf{制造流程概要}：集成电路的制造始于高纯度单晶硅棒，经
切片、研磨与抛光得到镜面般平整的晶圆。此后在晶圆上反复进行
氧化、光刻、刻蚀、掺杂（扩散或离子注入）与薄膜沉积等工序，
逐层构建出晶体管结构与多层金属互连。其中光刻是决定特征尺寸
的关键：把掩模版上的图形经缩小投影曝光到涂有光刻胶的晶圆
表面，再显影、刻蚀而成，可分辨的最小线宽受曝光波长与透镜
数值孔径限制。整片晶圆完成后先进行探针测试，剔除不合格的
裸片，再经划片、装架、键合与封装，最后进行老化筛选与成品
测试。由于工序繁多且每步都要求近乎无尘的环境，成品率成为
成本的决定因素：单片芯片面积越大，落入致命缺陷的概率越高，
成品率越低。这解释了为什么集成电路的价格与其面积并非线性
关系，也解释了为什么工艺进步在提高集成度的同时反而能降低
单位功能的成本。

\textbf{数字逻辑系列与主要参数}：通用数字集成电路按电路结构
形成若干互相兼容的系列。74 系列是应用最广的 TTL 及其兼容
产品，其中 74LS 为低功耗肖特基型，74HC 为高速 CMOS 型，
74HCT 则在 CMOS 工艺上保持与 TTL 兼容的输入电平；4000（CD4000）
系列为早期 CMOS 逻辑，电源电压范围宽（可达
$3\sim18\ \mathrm{V}$）、功耗极低，但速度较慢、驱动能力弱。
选用与互连时须核对四个参数。其一是电平规范，即输入高电平
下限 $V_{IH}$、输入低电平上限 $V_{IL}$ 与输出电平范围，
二者之差构成噪声容限，容限越大抗干扰能力越强。其二是传输
延迟时间 $t_{pd}$，它与工作频率上限直接相关。其三是扇出
系数，即一个输出端可驱动的同类输入端数目，受输出电流能力
限制；驱动继电器、发光二极管等大电流负载时必须外加驱动级。
其四是功耗，TTL 的静态功耗大而与频率无关，CMOS 的功耗则
主要随频率线性增长。不同系列混用时，还须注意电源电压与
电平的转换，必要时插入电平转换芯片或采用漏极开路输出加
上拉电阻的方式。

\textbf{通信设备中的典型集成电路}：在现代通信整机中，集成电路
几乎覆盖了全部信号链路。射频前端由低噪声放大器、混频器、
压控振荡器和锁相环频率合成器组成，其中锁相环通过比较参考频率
与压控振荡器分频后的相位来锁定输出频率，输出频率满足
$f_o = N f_{ref}$，$N$ 为可编程分频比，这是频率合成与信道选择
的基础。中频与基带部分由可变增益放大器、滤波器、ADC/DAC 和
数字信号处理器构成，负责解调、纠错与信源编解码。控制部分由
微控制器完成人机接口、频率与功率设置、状态监测与告警。此外，
现场可编程门阵列（FPGA）以其可重构性广泛用于协议处理与算法
验证，而把处理器核、存储器、外设接口与射频模块集成于单片的
片上系统（SoC），则使小型化手持终端成为可能。

\textbf{使用与维护要点}：集成电路的使用有若干必须遵守的规则。
第一是静电防护，CMOS 器件的栅极氧化层极薄，人体静电可达数千伏，
足以造成永久击穿或潜在损伤，故存取时须使用防静电包装、
接地手环与防静电烙铁。第二是电源去耦，每片集成电路的电源引脚
附近都应就近配置去耦电容，通常以 $0.1\ \mu$F 陶瓷电容旁路高频
噪声、以数 $\mu$F 至数十 $\mu$F 电容提供瞬态电流，二者并联可在
宽频带内维持低阻抗。第三是电源电压与极性，绝不可超出绝对最大
额定值，也不可反接；多电源芯片必须遵守规定的上电顺序。第四是
未用输入端的处理，CMOS 电路的悬空输入端极易受干扰而使电路进入
不定状态并产生大电流，必须按逻辑要求接至高电平或低电平。第五是
带电插拔的禁止，热插拔会在管脚上产生瞬态过压，除专门设计支持
热插拔的器件外一律应断电操作。

在检修方面，由于集成电路内部不可修复，故障判断只能依据外部
表现：先测量电源引脚电压是否正常、有无异常发热，再按信号流向
逐级测量输入输出电平与波形，用替换法最终确认。对贴装封装的
芯片，拆焊须使用热风枪并严格控制温度与时间，避免因热应力导致
基板起泡或相邻元件移位。理解集成电路的分级、工艺、封装与使用
规范，是掌握现代通信设备结构与维修方法的必要基础。
